%!TEX root = ../../main.tex
\section{규모 계급과 코호트 T, S}
\label{sec:eng-scale}

규모는 예측 시점 이후에 결정되는 사후 속성이며 보고와 코호트 구분에만 쓴다. 각 팀의 참여 인원은 교전의 킬에서 킬러·희생자·어시스트로 나타난 그 팀의 챔피언 수에, $\tcut$과 마지막 킬 사이에 앵커로부터 3{,}000 단위 안에서 발생한 다른 timeline 사건의 행위자를 더한 수이다. 계급은 작은 쪽의 인원 $\nmin$으로 정한다. 정의 상수 추정 코퍼스(\numEngLineage{} 교전)에서 $\nmin$은 1, 2, 3, 4, 5인 경우가 각각 \shareNminOne{}, \shareNminTwo{}, \shareNminThree{}, \shareNminFour{}, \shareNminFive{}\%로 작은 쪽 인원만으로는 antimode가 없고, 결합 분포의 대각선은 2v2(14.7\%)가 모드, 4v4(6.2\%)가 골, 5v5(7.2\%)가 두 번째 모드를 이룬다. 비지도 행동학이 밀도 봉우리에서 행동 상태를 읽듯 \cite{berman2014mapping}, 계급은 이 모드를 따른다. \emph{pick}은 $\nmin \le 1$(\numPickLineage{}, \sharePick{}\%), \emph{소규모 교전}(skirmish)은 $2 \le \nmin \le 3$(\numSkirmishLineage{}, \shareSkirmish{}\%), \emph{한타}(teamfight)는 $\nmin \ge 4$(\numTeamfightLineage{}, \shareTeamfight{}\%)이다.

두 모드는 두 종류의 다인 교전을 지지한다. 그 사이의 절단 위치는 자료가 정해 주는 값이 아니므로, 절단을 4에 두는 것은 팀 전체가 싸운다는 커뮤니티의 감각을 따른 관례이다. 같은 out-of-fold 예측을 다시 점수화하면 pick AUC에서 한타 AUC를 뺀 값은 절단 3에서 \cutDiffThree{}, 4에서 \cutDiffFour{}, 5에서 \cutDiffFive{}로 부호가 절단을 따른다. 따라서 규모 계급 사이의 절대 성능 차이는 정의의 효과와 분리되지 않으며, 본 논문은 두 코호트를 각각 보고한다(해석 범위는 제 \ref{sec:disc-scope} 절). 학회 이후의 분석에서 나왔던 ``예측 가능성은 규모와 함께 커진다''는 결론은 이 부호 반전과 사후 시간 특징의 누출이 확인되어 철회되었다(제 \ref{sec:rel-encounter} 절).

본 논문의 예측 표본은 두 코호트이다. \textbf{한타 T}는 동결 파이프라인에서 플래그 \texttt{cohort == 1}, 곧 저장된 v3.3 참여 인원에 대해 $\min(\texttt{cluster\_blue}, \texttt{cluster\_red}) \ge 4$이며 v3.3의 한타 계급(\texttt{fine == 2})과 모든 집합에서 정확히 일치한다. \textbf{소규모 교전 S}는 \texttt{cohort == 0 \& fine == 1}, 곧 $2 \le \nmin \le 3$이다. pick($\nmin \le 1$; 15.16 교전의 약 \sharePickTest{}\%)은 사전에 제외되었고 동결 파이프라인의 결과가 없다. \pendingauthor{pick 제외 사유 문장 확정 (A1)} 두 코호트는 각자의 $q$와 기준선, 각자의 주 대비를 가지며 모든 결과는 코호트 안에서 읽는다. 코호트 소속은 예측 시점 이후의 정보이다.

정의 상수 추정에 쓴 코퍼스(\numMatchesLineage{} 경기, \numEngLineage{} 교전, 그중 한타 \numTeamfightLineage{})와 본 논문의 예측 코퍼스(\numMatchesTotal{} 경기, 합동 T \numTPooled{})는 필터가 다른 별개의 집합이며 각각의 표에서 보고한다.
